\section{Реализация}

\subsection{Технологический стек}

При разработке библиотеки одними из требований были кроссплатформенность и высокая производительность. Для обеспечения максимальной производительности желателен достаточно низкоуровневый язык. Чтобы предоставить возможность использования библиотеки из любых других языков программирования, необходим API на языке C: его ключевое преимущество --- стабильный C ABI, который поддерживается практически любым языком через механизмы FFI (Foreign Function Interface). 

Исходя из всех этих требований, внутренняя реализация библиотеки выполнена на языке C++ стандарта C++17. Использование C++17 позволило применить современные идиомы: умные указатели (std::unique\_ptr, std::shared\_ptr) для управления памятью, std::optional для безопасной обработки состояний, std::map для кэширования, а также structured bindings для улучшения читаемости кода. При этом пользователю предоставляется исключительно C API, что гарантирует совместимость и простоту интеграции.

Для автоматизации сборки выбран CMake (минимальной версии 3.15). CMake позволяет генерировать проектные файлы для различных сред (Makefiles, Ninja, Visual Studio) и управлять зависимостями. В конфигурации предусмотрены опции для включения/отключения отладочного вывода, а также возможность сборки библиотеки без операций вставки и удаления, что дает прирост производительности, если вышеупомянутые операции не требуются.

Внешние зависимости библиотеки включают в себя:
\begin{itemize}
    \item \textbf{Алгоритмы сжатия} --- zlib, zstd, lz4, brotli. Библиотека не реализует собственные алгоритмы, а предоставляет адаптеры для существующих высокопроизводительных реализаций.
    \item \textbf{Тестирование} --- Google Test.
    \item \textbf{Бенчмаркинг} --- Google Benchmark.
\end{itemize}

\subsection{Документация}

API-документация генерируется с помощью Doxygen. Комментарии в заголовочном файле \texttt{compio.h} оформлены в формате, понятном Doxygen, что позволяет автоматически создавать HTML-версию справочного руководства. Так же настроены GitHub Actions для автоматической генерации документации, которую можно найти по адресу https://lovyagin.github.io/compio/.

\subsection{API для языка C}

Библиотека предоставляет низкоуровневый интерфейс на языке C, определённый в заголовочном файле \texttt{compio.h}. Все функции используют префикс \texttt{compio\_} и оперируют двумя непрозрачными (opaque) типами-дескрипторами: \texttt{compio\_archive} (открытый архив) и \texttt{compio\_file} (открытый файл внутри архива). Такой подход обеспечивает строгую инкапсуляцию и позволяет использовать библиотеку в проектах на чистом C, а также в любых других языках, поддерживающих C-совместимые интерфейсы.

\paragraph{Открытие и закрытие архива.}
Для работы с архивом вызывается функция
\begin{verbatim}
compio_archive *compio_open_archive(const char *fp, const char *mode,
                                    const compio_config *c);
\end{verbatim}
Параметр \texttt{fp} задаёт путь к файлу архива, \texttt{mode} определяет режим доступа:
\begin{itemize}
    \item \texttt{"r"} --- открыть существующий архив только для чтения.
    \item \texttt{"r+"} --- открыть существующий архив для чтения и записи.
    \item \texttt{"w"} --- создать новый архив (если файл существует, он перезаписывается).
    \item \texttt{"w+"} --- создать новый архив для чтения и записи.
    \item \texttt{"a"} --- открыть существующий архив для добавления; все вновь открываемые файлы изначально позиционируются в конец, существующие данные сохраняются.
    \item \texttt{"a+"} --- эквивалентно \texttt{"a"}.
\end{itemize}
Третий аргумент должен быть указателем на структуру конфигурации \texttt{compio\_config}. При успешном открытии возвращается указатель на \texttt{compio\_archive}, при ошибке --- \texttt{NULL}, а код ошибки помещается в \texttt{errno}. Архив закрывается вызовом \texttt{compio\_close\_archive}, который сбрасывает все кэши, синхронизирует метаданные с диском и освобождает ресурсы.

\paragraph{Работа с файлами внутри архива.}
Файл внутри архива открывается функцией
\begin{verbatim}
compio_file *compio_open_file(const char *name, compio_archive *archive);
\end{verbatim}
Имя файла не должно превышать \texttt{COMPIO\_FNAME\_MAX\_SIZE} (32 символа по умолчанию). Возвращаемый дескриптор \texttt{compio\_file} используется для всех последующих операций чтения, записи и позиционирования. После завершения работы файл закрывается через \texttt{compio\_close\_file}.

\paragraph{Операции ввода-вывода.}
Основные функции для работы с данными:
\begin{itemize}
    \item \texttt{uint64\_t compio\_read(void *ptr, uint64\_t size, compio\_file *file)} --- читает до \texttt{size} байт из текущей позиции файла в буфер \texttt{ptr}, возвращает количество реально прочитанных байт.
    \item \texttt{uint64\_t compio\_write(const void *ptr, uint64\_t size, compio\_file *file)} --- записывает \texttt{size} байт из буфера \texttt{ptr} в текущую позицию файла, возвращает число записанных байт.
    \item \texttt{uint64\_t compio\_insert(const void *ptr, uint64\_t size, compio\_file *file)} --- вставляет \texttt{size} байт из буфера \texttt{ptr} в текущую позицию файла, сдвигая последующие байты вправо.
    \item \texttt{uint64\_t compio\_erase(uint64\_t size, compio\_file *file)} --- удаляет \texttt{size} байт с текущей позиции файла, сдвигая хвост файла влево.
    \item \texttt{int compio\_seek(compio\_file *file, int64\_t offset, uint8\_t origin)} --- перемещает текущую позицию на \texttt{offset} байт относительно текущей позиции в файле, начала файла или конца файла (управляется аргументом \texttt{origin}).
    \item \texttt{uint64\_t compio\_tell(compio\_file *file)} --- возвращает текущую позицию в файле.
\end{itemize}
Все функции в случае ошибки устанавливают \texttt{errno} (например, \texttt{ENAMETOOLONG}, \texttt{ENFILE}, \texttt{EROFS}, \texttt{EINVAL}).

\paragraph{Конфигурация.}
Структура \texttt{compio\_config} объединяет все настраиваемые параметры: алгоритм сжатия (\texttt{compressor}), степень B-дерева (\texttt{tree\_degree}), целевой размер блока (\texttt{block\_size}) и границы минимального/максимального размера (\texttt{block\_size\_\_minimum} и \texttt{block\_size\_\_maximum}), ёмкость кэшей для узлов и блоков данных (\texttt{cache\_size\_\_nodes} и \texttt{cache\_size\_\_blocks}). Для удобства предоставлена функция \texttt{compio\_build\_default\_config}, заполняющая структуру значениями по умолчанию.

Поля \texttt{tree\_degree} и \texttt{block\_size} пользователь может задать равными нулю и тогда при открытии существующего архива они будут прочитаны из заголовка. Это поведение работает только в случае, когда открывается уже существующий архив, а не создается новый.

\paragraph{Выбор алгоритма сжатия.}
Для каждого встроенного алгоритма реализована своя функция-конструктор: \texttt{compio\_build\_zlib\_compressor}, \texttt{compio\_build\_lz4\_compressor}, \texttt{compio\_build\_zstd\_compressor}, \texttt{compio\_build\_brotli\_compressor}, а также \texttt{compio\_build\_dummy\_compressor} (отсутствие сжатия). Чтобы использовать собственный алгоритм, пользователь должен самостоятельно заполнить структуру \texttt{compio\_compressor}, указав функции \texttt{compress}, \texttt{decompress}, \texttt{get\_buf\_size} и установив поле \texttt{compression\_type} в \texttt{COMPIO\_COMPRESS\_CUSTOM}.

Созданный экземпляр \texttt{compio\_compressor} помещается в поле \texttt{compressor} структуры \texttt{compio\_config}. Затем эта конфигурация передаётся в функцию \texttt{compio\_open\_archive} при создании нового архива (режимы \texttt{"w"}, \texttt{"w+"} или \texttt{"a"}). Чтобы файл-контейнер был самодостаточным, и его можно было арспаковать без дополнительной информации про алгоритм сжатия, библиотека сохраняет значение \texttt{compression\_type} из компрессора в заголовке контейнера.

Перед открытием существующего архива пользователь может определить тип сжатия, вызвав функцию \texttt{compio\_get\_compression\_type}. Затем он обязан самостоятельно инициализировать компрессор, соответствующий этому типу. Для встроенных алгоритмов это можно сделать с помощью вспомогательной функции \texttt{compio\_build\_compressor\_by\_type}. Для типа \texttt{COMPIO\_COMPRESS\_CUSTOM} пользователь должен предоставить собственную реализацию компрессора, идентичную той, что использовалась при создании архива. 

Библиотека не создаёт компрессор автоматически и не изменяет поле \texttt{compressor} в переданной конфигурации. Если тип сжатия, прочитанный из заголовка архива, не соответствует типу, указанному в предоставленном пользователем компрессоре (например, архив требует ZLIB, а передан LZ4), или для кастомного типа передан некорректный компрессор, функция \texttt{compio\_open\_archive} вернёт ошибку. Такой механизм даёт пользователю полный контроль над процессом, но требует явных действий для согласования компрессора с архивом.

Внутри библиотеки все операции сжатия и разжатия выполняются исключительно через вызовы функций, на которые указывают поля \texttt{compress}, \texttt{decompress} и \texttt{get\_buf\_size}. Это позволяет абстрагироваться от конкретного алгоритма и использовать любые кодеки, в том числе и предоставленные пользователем.

\paragraph{Пример использования.}
Полный пример работы с библиотекой на языке C:
\begin{verbatim}
#include "compio.h"

int main() {
    compio_config cfg;
    compio_build_default_config(&cfg);
    compio_build_lz4_compressor(&cfg.compressor);

    compio_archive *arch = compio_open_archive("data.bin", "w+", &cfg);
    compio_file *f = compio_open_file("A", arch);

    const char *data = "Hello, world!";
    compio_write(data, 14, f);

    char buf[14];
    compio_seek(f, 0, COMPIO_SEEK_SET);
    compio_read(buf, 14, f);

    compio_close_file(f);
    compio_close_archive(arch);

    return 0;
}
\end{verbatim}